\NativeSlot{43}{74}{%
% p011_s01

\ExamQuestion{10.}{5}%
工程队在修建高速路时，有时需要将弯曲的道路改直来缩短路程，这样做的道理是\mbox{（\hspace{1.8em}）}。\par
\ExamOptions{A. 两点确定一条直线}{B. 两点之间线段最短}{C. 过一点有无数条直线}{D. 不确定}
}
\NativeSlot{167}{87}{%
\ExamGroup{三、计算。{\normalfont\kaishu\fontsize{10.8bp}{14bp}\selectfont （共24分）}}

\noindent \textbf{1.}\quad 列竖式计算。（12分）\par\smallskip
\noindent%
% p011_s02
\makebox[.249\linewidth][l]{\(114\times23=\hspace{1.8em}\)}%
% p011_s03
\makebox[.249\linewidth][l]{\(448\div48=\hspace{1.8em}\)}%
% p011_s04
\makebox[.249\linewidth][l]{\(372\div62=\hspace{1.8em}\)}%
% p011_s05
\makebox[.249\linewidth][l]{\(598\div18=\hspace{1.8em}\)}\par
}
\NativeSlot{316}{69}{%
\noindent \textbf{2.}\quad 计算下列各题。（12分）\par\smallskip
\par\noindent%
\ExamTile{.323}{%
% p011_s06

\noindent \(1008\div[(73+32)\div5]\)\par

\ExamWork{3}
}\hfill%
\ExamTile{.323}{%
% p011_s07

\noindent \(966\div(27+19)\times45\)\par

\ExamWork{3}
}\hfill%
\ExamTile{.323}{%
% p011_s08

\noindent \(28\div4+63\times25\)\par

\ExamWork{3}
}\par
}
\NativeSlot{443}{103}{%
\ExamGroup{四、画一画。{\normalfont\kaishu\fontsize{10.8bp}{14bp}\selectfont （共11分）}}

% p011_s09

\ExamQuestion{1.}{5}%
过\(A\)点画出已知直线\(l\)的平行线和垂线。\par
\BookFigure{figures/p011_s09-01.tex}
}
\NativeSlot{621}{104}{%
% p011_s10

\ExamQuestion{2.}{5}%
补充两条边，使下面的图形变成一个长方形。\par
\BookFigure{figures/p011_s10-01.tex}
}
